\documentclass[UTF8,12pt]{ctexart}
\usepackage[a4paper,margin=24mm]{geometry}
\usepackage{amsmath,amssymb,bm,graphicx,adjustbox}
\newcommand{\fitmath}[1]{\begin{adjustbox}{max width=\linewidth}$\displaystyle #1$\end{adjustbox}}
\setlength{\parindent}{0pt}
\setlength{\parskip}{.7em}
\sloppy
\allowdisplaybreaks
\begin{document}
\section*{1996 年考研数学三 · 参考答案}
\subsection*{1996 年真题参考答案（试卷 IV）}

\subsection*{一、填空题}

（1）\(\displaystyle \dfrac{\displaystyle dx}{\displaystyle x(1+\ln y)}\)。 （2）\(\displaystyle -\dfrac{\displaystyle 1}{\displaystyle 3}\sqrt{(1-x^2)^3}+C\)，其中 \(\displaystyle C\) 为任意常数。 （3）\(\displaystyle \dfrac{\displaystyle c}{\displaystyle a}\geq0\)，\(\displaystyle b\) 任意。 （4）\(\displaystyle (1,\allowbreak 0,\allowbreak \ldots,\allowbreak 0)^{\mathrm T}\)。 （5）\(\displaystyle (4.412,\allowbreak 5.588)\)。


\subsection*{二、选择题}

（1）D。 （2）A。 （3）C。 （4）D。 （5）B。


三、（1）\(\displaystyle f^{\prime}(x)=\begin{cases}\dfrac{\displaystyle xg^{\prime}(x)-g(x)+(x+1)e^{-x}}{\displaystyle x^2},\allowbreak &x\ne0,\allowbreak \\\dfrac{\displaystyle g^{\prime\prime}(0)-1}{\displaystyle 2},\allowbreak &x=0.\end{cases}\)（2）\(\displaystyle f^{\prime}(x)\) 在 \(\displaystyle (-\infty,\allowbreak +\infty)\) 上连续。


四、0。


五、\(\displaystyle \ln2\)。


六、证明略。（可利用积分中值定理和罗尔定理。）


七、（1）当 \(\displaystyle 0<p<\sqrt{\dfrac{\displaystyle b}{\displaystyle c}}(\sqrt a-\sqrt{bc})\) 时，随单价 \(\displaystyle p\) 增加，相应的销售额增加；当 \(\displaystyle p>\sqrt{\dfrac{\displaystyle b}{\displaystyle c}}(\sqrt a-\sqrt{bc})\) 时，随单价 \(\displaystyle p\) 增加，相应的销售额减少。


（2）当 \(\displaystyle p=\sqrt{\dfrac{\displaystyle b}{\displaystyle c}}(\sqrt a-\sqrt{bc})\) 时，销售额取最大值，最大销售额是 \(\displaystyle (\sqrt a-\sqrt{bc})^2\)。


八、原方程的通解为 \(\displaystyle y+\sqrt{x^2+y^2}=C\)（\(\displaystyle x>0\)），其中 \(\displaystyle C\) 为任意非零常数。当 \(\displaystyle x<0\) 时，原方程的解与 \(\displaystyle x>0\) 时相同。


九、（1）\(\displaystyle y=2\)。（2）\(\displaystyle P=\begin{pmatrix}1&0&0&0\\0&1&0&0\\0&0&1&-\dfrac{\displaystyle 4}{\displaystyle 5}\\0&0&0&1\end{pmatrix}\)，\(\displaystyle (AP)^{\mathrm T}(AP)\) 为对角矩阵。


十、证明略。（可利用向量组线性无关的定义。）


十一、5.216 万元。


十二、\(\displaystyle p=\dfrac{\displaystyle 19}{\displaystyle 36},\allowbreak \ q=\dfrac{\displaystyle 1}{\displaystyle 18}\)。


十三、证明略。（使用中心极限定理。）\(\displaystyle Z_n\) 近似服从参数为 \(\displaystyle \mu=\alpha_2,\allowbreak \ \sigma^2=\dfrac{\displaystyle \alpha_4-\alpha_2^2}{\displaystyle n}\) 的正态分布。


\subsection*{1996 年真题参考答案（试卷 V）}

\subsection*{一、填空题}

（1）【同试卷 IV 第一、（1）题】 （2）【同试卷 IV 第一、（2）题】 （3）\(\displaystyle \dfrac{\displaystyle 5}{\displaystyle 32}\)。 （4）\(\displaystyle 1-a+a^2-a^3+a^4-a^5\)。 （5）\(\displaystyle \dfrac{\displaystyle 11}{\displaystyle 24}\)。


\subsection*{二、选择题}

（1）D。 （2）A。 （3）【同试卷 IV 第二、（3）题】 （4）【同试卷 IV 第二、（4）题】 （5）B。


三、【同试卷 IV 第三题】


四、\(\displaystyle -2e^{-x^2y^2}\)。


五、【同试卷 IV 第五题】


六、【同试卷 IV 第六题】


七、（1）证明略。（利用定积分求面积。）（2）记抛物线与两坐标轴所围图形绕 x 轴旋转所得旋转体体积为 \(\displaystyle V_1\)，抛物线与 x 轴所围图形绕 x 轴旋转所得旋转体体积为 \(\displaystyle V_2\)，则 \(\displaystyle V_1/V_2=19/8\)。


八、证明略。（可利用积分中值定理和罗尔定理。）


九、当 \(\displaystyle t\ne-2\) 时，方程组无解；当 \(\displaystyle t=-2\) 时，方程组有解，此时分两种情况：当 \(\displaystyle p=-8\) 时，方程组的通解为 \(\displaystyle x=k_1(4,\allowbreak -2,\allowbreak 1,\allowbreak 0)^{\mathrm T}+k_2(-1,\allowbreak -2,\allowbreak 0,\allowbreak 1)^{\mathrm T}+(-1,\allowbreak 1,\allowbreak 0,\allowbreak 0)^{\mathrm T}\)，其中 \(\displaystyle k_1,\allowbreak k_2\) 为任意常数；当 \(\displaystyle p\ne-8\) 时，方程组的通解为 \(\displaystyle x=k(-1,\allowbreak -2,\allowbreak 0,\allowbreak 1)^{\mathrm T}+(-1,\allowbreak 1,\allowbreak 0,\allowbreak 0)^{\mathrm T}\)，其中 \(\displaystyle k\) 为任意常数。


十、\(\displaystyle \dfrac{\displaystyle 4}{\displaystyle 3}\)。


十一、【同试卷 IV 第十一题】


十二、\(\displaystyle G(t)=\begin{cases}1-e^{-3\lambda t},\allowbreak &t>0,\allowbreak \\0,\allowbreak &t\leq0.\end{cases}\)\(\displaystyle T\) 服从参数为 \(\displaystyle 3\lambda\) 的指数分布。

\end{document}
